\subsection*{Rod Geometry and Kinematic Parameterization}

We model the human spine as a continuous elastic Cosserat rod, defined by a centerline curve $\mathbf{r}(s,t) \in \mathbb{R}^3$ parameterized by arc length $s \in [0, L]$. To each point $s$, we attach an orthonormal material frame $\{\mathbf{d}_1(s,t), \mathbf{d}_2(s,t), \mathbf{d}_3(s,t)\}$ describing the cross-sectional orientation~\cite{goriely2017mathematics, gazzola2018forward}. The director $\mathbf{d}_3(s,t)$ aligns with the tangent to the centerline in the absence of shear strain, while $\mathbf{d}_1(s,t)$ and $\mathbf{d}_2(s,t)$ define the principal axes of the vertebral cross-section. The kinematics are governed by the evolution equations:

\begin{align}
    \mathbf{r}'(s) &= \mathbf{v}(s) = \mathbf{d}_3(s) + \boldsymbol{\varepsilon}(s), \\
    \mathbf{d}_i'(s) &= \mathbf{u}(s) \times \mathbf{d}_i(s),
    \label{eq:cosserat_kinematics}
\end{align}

where $\mathbf{v}(s)$ is the tangent vector and $\boldsymbol{\varepsilon}(s)$ denotes the strain vector (shear and extension). The curvature vector $\boldsymbol{\kappa}(s)$ and twist $\tau(s)$ are encoded in the Darboux vector $\mathbf{u}(s) = \kappa_1 \mathbf{d}_1 + \kappa_2 \mathbf{d}_2 + \tau \mathbf{d}_3$. We assign the physical boundary conditions: $s=0$ corresponds to the sacrum (clamped: $\mathbf{r}(0)=\mathbf{0}, \mathbf{d}_i(0)=\mathbf{e}_i$) and $s=L$ to the cranium (free), defining the mechanical axis of the column.

\subsection*{Developmental Information Fields from Morphogenetic Patterning}

We formalize the developmental input by defining a time-dependent scalar information field $I(s, t) \in [0, 1]$ representing the coarse-grained expression level of anterior-posterior patterning morphogens (e.g., HOX genes). Following the bimodal Gaussian parameterization (Methods Eq.~\ref{eq:info_field_spinal}), we model $I(s)$ as a bimodal Gaussian distribution with peaks centered at the cervical ($s_c=0.80$) and lumbar ($s_l=0.25$) enlargements, with normalized amplitudes $A_c=0.5$ and $A_l=0.7$, respectively. These peaks correspond to regions of high vertebral identity specification. The spatial gradient $\nabla I$ serves as a ``curvature generator'' via the geometric coupling $\chi_\kappa$:
\begin{equation}
    \boldsymbol{\kappa}_{\mathrm{rest}}(s) = \boldsymbol{\kappa}_{\mathrm{gen}} + \chi_\kappa \nabla I(s).
    \label{eq:kappa_rest}
\end{equation}

Simultaneously, the field exerts a stiffness via the morphomechanical coupling $\chi_M$, generating the active biological moment $\mathbf{M}_{bio}$ to oppose gravity:

\begin{equation}
    \mathbf{M}_{bio} = \chi_M (\mathbf{\Lambda} \cdot \nabla I).
    \label{eq:active_moment}
\end{equation}

Stability is quantified by the Bio-Gravitational Number $\mathcal{B}_g$:

\begin{equation}
    \mathcal{B}_g = \frac{\chi_M \langle |\nabla I| \rangle}{\rho A g L^2}.
    \label{eq:bg_number}
\end{equation}

When $\mathcal{B}_g > 1$, the organism dominates gravity. This framework is physically grounded in the structural specificity of HOX proteins (e.g., HOXC8, HOXB13). AlphaFold structural analysis confirms that these factors possess rigid DNA-binding domains (e.g., HOXC8, UniProt P31273), which are capable of enforcing sharp identity boundaries.
